% Nguồn: docs/project/brokerguard_v3_formal_model.md (mục 1--6). Bằng chứng trong Bảng 2: facts.md mục Test-4 (ablation
% theo cơ chế, bằng chứng cơ chế mức ca, hai họ M5 còn lại) và mục "Bổ sung cho bài v3" (C1 offline trên Dev-val).
% Mệnh đề 3 chỉ bảo đảm literal ghim CÓ MẶT trong data(a), không bảo đảm không có giá trị khác bên cạnh.
% Mã giữ lại ở mục này: (G1)--(G4) và (A1)--(A5), vì chứng minh và Bảng 2 tham chiếu tới chúng.
% 20/11 (lượt 3): G(u), Sigma định nghĩa tại đây; vị từ gây hại gọi bằng tên (chuyển hướng, rò rỉ...) thay vì H_dest...; pins bằng lời.
\section{Mô hình hình thức}\label{sec:formal}

Mục này xác định khi nào các kiểm tra ở Mục~\ref{sec:method} đủ ngăn các kiểu tác hại ở Mục~\ref{sec:threat}. Gọi
$G(u)$ là tập cặp công cụ và đích mà người dùng thực sự cho phép, và $\Sigma$ là tập giá trị mật mà người dùng không cho
tiết lộ. Trạng thái của broker sau bước $t$ là $s_t = (u, P(u), C_t, H_t)$, với $H_t$ là lịch sử lời gọi đã thực thi và
kết quả. $P(u)$ được dựng một lần trước bước đầu. Gọi $R$ là tập công cụ đọc và $S$ là tập sink có trường đích. Cổng xét
mỗi lời gọi đề xuất $a_t$ và trả về cho qua, từ chối hoặc cần xác nhận, trong đó cần xác nhận được xử lý như từ chối. Lời
gọi đọc luôn được cho qua. Với $a_t \in S$, cổng cho qua khi và chỉ khi (G1) \rSig{} không chặn $a_t$, (G2) $a_t$ thỏa
\eqref{eq:c1}, (G3) có ít nhất một quyền $g$ thỏa (G2) mà mọi literal được ghim của $g$ đều có mặt trong $\data(a_t)$ và
không boolean nào bị đảo, và (G4) điều kiện \eqref{eq:c2} sai. Gọi $\hat G(u)$ là tập cặp công cụ và đích thỏa (G2). Khi
cổng cho qua, broker thực thi $a_t$, ghi lời gọi và kết quả $o_t$ vào $H_t$, rồi cập nhật $C_{t+1} = C_t \cup L(o_t)$
với $L$ là hàm gắn nhãn kết hợp \rFlow{} và \pTwo. Câu trả lời cuối được phát hành sau khi che mọi giá trị của $C_t
\setminus V(u)$. Năm giả định sau giới hạn phạm vi của các mệnh đề.
\begin{enumerate}[label=(A\arabic*), leftmargin=*]
\item \emph{Trung gian toàn phần.} Mọi lời gọi công cụ đều đi qua broker, và mọi công cụ có tác dụng phụ đều thuộc $S$.
\item \emph{Bộ trích đúng.} $\hat G(u) \subseteq G(u)$. Chiều ngược lại chỉ ảnh hưởng tới tiện ích.
\item \emph{Độ phủ nhãn.} Mọi giá trị của $\Sigma$ xuất hiện trong một kết quả $o_t$ đều thuộc $C_{t+1}$.
\item \emph{Lớp biến đổi.} Nếu tác tử làm lộ $v$ trong một chuỗi $w$ thì $v \sub w$. Mã hóa, viết số bằng chữ hay chia
  một giá trị qua nhiều lời gọi nằm ngoài lớp này.
\item \emph{Literal ghim đủ.} Với quyền $g$ cấp cho lời gọi $a$, các literal được ghim của $g$ gồm mọi giá trị và
  boolean mà người dùng đã nêu cho $a$.
\end{enumerate}

\begin{proposition}[An toàn cấp quyền]\label{prop:auth}
Dưới (A1) và (A2), mọi lời gọi sink được thực thi có cặp công cụ và đích thuộc $G(u)$. Vì vậy không quỹ đạo nào gây hại
kiểu chuyển hướng, và nhiệm vụ chỉ cần đọc không gây tác dụng phụ.
\end{proposition}
\begin{proof}
Theo (A1), một sink $a$ chỉ được thực thi khi cổng cho qua, nên (G2) đúng và cặp công cụ, đích của $a$ thuộc $\hat G(u)$,
mà $\hat G(u) \subseteq G(u)$ theo (A2). Nhiệm vụ chỉ cần đọc có $G(u)$ rỗng, nên không sink nào được thực thi. \qed
\end{proof}

\begin{proposition}[Bảo mật]\label{prop:conf}
Dưới (A1), (A3) và (A4), không giá trị nào của $\Sigma$ xuất hiện trong $\data(a)$ của một sink được thực thi hay trong
câu trả lời cuối đã phát hành. Vì vậy không quỹ đạo nào gây hại kiểu rò rỉ ở các dạng thuộc (A4).
\end{proposition}
\begin{proof}
Giả sử $v \in \Sigma$ xuất hiện trong $\data(a)$ của một sink $a$ được thực thi ở bước $t$. Vì $v \notin V(u)$, tác tử
chỉ có được $v$ qua một kết quả trước bước $t$, nên theo (A3) $v \in C_t$. Theo (A4), $v \sub \data(a)$, nên
\eqref{eq:c2} đúng và (G4) sai, mâu thuẫn với (A1). Với câu trả lời cuối, phép che thay mọi dạng bề mặt khớp
$\kappa(v)$ trước khi phát hành, nên theo (A4) không còn dạng nào của $v$. \qed
\end{proof}

\begin{proposition}[Toàn vẹn]\label{prop:int}
Dưới (A1), (A2) và (A5), mọi lời gọi sink được thực thi chứa mọi giá trị người dùng đã nêu cho lời gọi đó và không mang
boolean ngược với boolean người dùng đã nêu.
\end{proposition}
\begin{proof}
Theo (A1), sink $a$ chỉ được thực thi khi (G2) và (G3) đúng, tức có quyền $g$ cấp cho $a$ mà mọi literal được ghim của
$g$ có mặt trong $\data(a)$ và không boolean nào bị đảo. Theo (A2), $g$ ứng với một quyền người dùng thật sự cấp, và theo
(A5), các literal được ghim của $g$ gồm mọi giá trị người dùng đã nêu cho $a$. \qed
\end{proof}

\noindent Ba mệnh đề nêu điều kiện đủ dưới các giả định, không phải bảo đảm vô điều kiện. Mệnh đề~\ref{prop:int} không
loại trừ việc một giá trị khác được thêm cạnh literal đã ghim. Lời gọi đọc và sink không có trường đích nằm ngoài $S$, nên
một công cụ đọc nhận URL do mô hình chọn vẫn có thể thành kênh rò rỉ. Bảng~\ref{tab:proven} phân biệt thuộc tính được
chứng minh với thuộc tính chỉ kiểm chứng bằng thực nghiệm, và các giả định thực nghiệm được kiểm tra ở
Mục~\ref{sec:results}.

% Bảng 2. Nguồn: facts.md mục Test-4 (H3 ablation; họ bị hại theo cơ chế: −C1 thêm 1 ER và 1 PT, −C2 thêm 2 SL và 2 AL;
% lộ câu cuối v3 0/720; hai họ M5 còn lại; nguyên văn = biến đổi) và mục "Bổ sung cho bài v3" (C1 offline Dev-val:
% tiện ích 155/155 và 65/65, M6 35/35, v1 15). (A1): tests/unit/test_brokerguard_v3_v1.py. Bọc adjustbox max width.
\begin{table}[!htbp]
\centering
\caption{Thuộc tính được chứng minh và thuộc tính chỉ được đo bằng thực nghiệm. Số trên tập phát triển được dùng khi
thiết kế, số trên tập kiểm thử có sau khi đóng băng. \cOne{} \pOne, \cTwo{} \pTwo, \cThree{} \pThree.}\label{tab:proven}
\footnotesize
\setlength{\tabcolsep}{3pt}
\begin{adjustbox}{max width=\textwidth}
\begin{tabularx}{\textwidth}{@{}>{\raggedright\arraybackslash}p{2.75cm}>{\raggedright\arraybackslash}p{2.2cm}%
>{\raggedright\arraybackslash}X@{}}
\toprule
Thuộc tính & Trạng thái & Bằng chứng \\
\midrule
Mệnh đề~\ref{prop:auth}--\ref{prop:int} & Chứng minh dưới giả định & Mục~\ref{sec:formal} \\
(A1) trung gian toàn phần & Theo cấu trúc mã, có kiểm thử & Broker là đường duy nhất tới công cụ. Công cụ đọc và sink
  không có đích nằm ngoài $S$ \\
(A2) bộ trích đúng & Thực nghiệm & Tập phát triển: mọi sink lành tính được cho phép, chặn đủ 35 sink tác dụng phụ ngoài
  yêu cầu (\rDest: 15). Tập kiểm thử: bỏ \cOne{} làm lọt thêm 2 họ \\
(A3) độ phủ nhãn & Thực nghiệm, gián tiếp & Tập kiểm thử: bỏ \cTwo{} làm lọt thêm 4 họ rò rỉ. \ours{} không để lộ giá
  trị nào ở câu trả lời cuối \\
(A4) lớp biến đổi & Giới hạn đã nêu & Bộ chấm biến đổi cho cùng kết quả với so khớp nguyên văn ở mọi cấu hình, nên giả
  định chưa bị thử thách \\
(A5) literal ghim đủ & Thực nghiệm, chưa đủ & Tập kiểm thử: 2 họ sửa giá trị số hoặc đơn vị lọt qua cả \strSand{} và
  \ours. Bỏ \cThree{} không làm lọt thêm họ nào \\
Lời nhắc, \rSig & Chỉ thực nghiệm & Không có bảo đảm. Đo qua ablation và tấn công thích nghi (Mục~\ref{sec:results})
  \\
\bottomrule
\end{tabularx}
\end{adjustbox}
\end{table}

